% !TeX root = booklet-02.tex
\chapter*{تمرین نشتی و جمع‌بندی مسیرهای جریان}
\addcontentsline{toc}{chapter}{تمرین نشتی و جمع‌بندی مسیرهای جریان}
\section*{تمرین \secnum{۲.۱۰}، اثر کاهش آستانه بر جریان روشن و نشتی}
\paragraph*{صورت مسئله و شرایط ثابت در مقایسه}
ترانزیستور \lr{nMOS} با \(V_{DD}=\qty{1.2}{\volt}\) و
\(V_t=\qty{0.4}{\volt}\) داده شده است. پیشنهاد طراحی، کاهش آستانه
به اندازه \(\qty{100}{\milli\volt}\) است. ضریب تغییر جریان اشباع
ایده‌آل در \(V_{GS}=V_{DS}=V_{DD}\)، ضریب تغییر نشتی زیرآستانه در
گیت صفر و دمای اتاق با \(n=1.4\)، و همان ضریب نشتی در
\(\qty{120}{\degreeCelsius}\) خواسته شده‌اند.
طبق صورت سؤال، آستانه مستقل از دما فرض می‌شود
\cite[تمرین \secnum{۲.۱۰}، صفحه ۹۶]{weste}.

کاهش مثبت آستانه را به صورت زیر تعریف می‌کنیم.
\[
 \Delta V_t=V_{t,\mathrm{old}}-V_{t,\mathrm{new}}
 =\qty{0.100}{\volt}.
\]
پس \(V_{t,\mathrm{new}}=\qty{0.3}{\volt}\) است.
در هر مقایسه، بایاس، هندسه و ضرایب مدل جز آستانه یکسان‌اند.
در دو بخش نشتی، نسبت جریان جدید به قدیم در دمای همان بخش محاسبه
می‌شود.

\paragraph*{بخش الف، جریان اشباع ایده‌آل}
مدل \lr{Shockley} در اشباع می‌دهد
\(I_{\mathrm{sat}}=\beta(V_{GS}-V_t)^2/2\).
اضافه‌ولتاژ قدیم و جدید به‌ترتیب
\(\qty{0.8}{\volt}\) و \(\qty{0.9}{\volt}\) هستند.
چون \(V_{DS}=\qty{1.2}{\volt}\) از هر دو اضافه‌ولتاژ بزرگ‌تر
است، شرط اشباع هر دو حالت برقرار است. تقسیم دو جریان، ضریب
\(\beta/2\) را حذف می‌کند و نتیجه می‌دهد
\begin{equation}
 r_{\mathrm{on}}=
 \frac{I_{\mathrm{sat,new}}}{I_{\mathrm{sat,old}}}
 =\left(\frac{1.2-0.3}{1.2-0.4}\right)^2
 =\left(\frac{0.9}{0.8}\right)^2
 =\frac{81}{64}=1.265625\approx\boxed{1.27}.
 \label{eq:leakage-ex10-on-ratio}
\end{equation}
جریان حدود \num{26.6} درصد افزایش می‌یابد. تبدیل این نسبت به
تأخیر مدار به بار و جریان متوسط گذار نیز وابسته است.
اگر آن جریان متوسط با همین نسبت تغییر کند و بار و دامنه ثابت باشند،
برآورد \(t\propto1/I_{\mathrm{avg}}\)، نسبت تأخیر
\(64/81\approx0.790\) می‌دهد. این برآورد شرطی از مدل جریان
متوسط استفاده می‌کند.

\paragraph*{بخش ب، نسبت نشتی در دمای اتاق}
در رابطه \eqref{eq:leakage-subthreshold-model}، صورت توان حالت
جدید را از صورت توان حالت قدیم کم می‌کنیم. جمله‌های گیت، بدنه و
درین برابرند و تفاوت برابر \(\Delta V_t\) می‌شود.
با \(I_{ds0}\)، \(n\) و عامل درین یکسان در همان دما، داریم
\begin{align}
 r_{\mathrm{sub}}(T)
 &=\frac{I_{ds0}e^{(V_{GS}-V_{t,\mathrm{new}}+\eta V_{DS}-k_\gamma V_{SB})/(nv_T)}F(V_{DS})}
 {I_{ds0}e^{(V_{GS}-V_{t,\mathrm{old}}+\eta V_{DS}-k_\gamma V_{SB})/(nv_T)}F(V_{DS})}
 \notag\\
 &=\boxed{\exp\!\left(\frac{\Delta V_t}{nv_T(T)}\right)}.
 \label{eq:leakage-ex10-ratio-derivation}
\end{align}
در این نسبت، اثرهای بدنه و \lr{DIBL} با فرض یکسان‌بودن در دو
نمونه حذف شده‌اند. \(V_{DS}=V_{DD}>0\) است، پس جریان مرجع مثبت
و تقسیم مجاز است.
با تقریب اتاق کتاب \(v_T=\qty{0.026}{\volt}\)، ابتدا توان را
بی‌بعد حساب می‌کنیم
\begin{equation}
 x=\frac{0.100}{1.4(0.026)}\approx2.7473,\qquad
 r_{\mathrm{sub,room}}=e^x\approx\boxed{15.6}.
 \label{eq:leakage-ex10-room}
\end{equation}
برای کنترل مستقل،
\(S=1.4(0.026)\ln10\approx\qty{83.81}{\milli\volt}\) به ازای
هر دهه است. کاهش \(\qty{100}{\milli\volt}\) آستانه برابر
\(100/83.81\approx1.193\) دهه است و
\(10^{1.193}\approx15.6\) همان نسبت را می‌دهد.

\paragraph*{بخش ج، نسبت نشتی در دمای بالاتر}
دمای مطلق از \(T[\mathrm{K}]=\theta[{}^\circ\mathrm{C}]+273.15\)
به دست می‌آید. نماد \(\theta\) دمای سلسیوس و \(T\) دمای کلوین
است. پس \(\qty{120}{\degreeCelsius}=\qty{393.15}{\kelvin}\).
با \(k=\qty{1.380649e-23}{\joule\per\kelvin}\) و
\(q=\qty{1.602176634e-19}{\coulomb}\)، ولتاژ حرارتی برابر است با
\begin{equation}
 v_T=\frac{kT}{q}\approx\qty{0.033879}{\volt}
 \approx\qty{34}{\milli\volt}.
 \label{eq:leakage-ex10-hot-thermal}
\end{equation}
با تقریب \(\qty{34}{\milli\volt}\)، نسبت در دمای بالا چنین است
\begin{equation}
 r_{\mathrm{sub,hot}}=
 \exp\!\left(\frac{0.100}{1.4(0.034)}\right)
 \approx8.173\approx\boxed{8.2}.
 \label{eq:leakage-ex10-hot}
\end{equation}
\(I_{ds0}(T)\) می‌تواند تابع دما باشد، اما در نسبت دو آستانه در
همان دما، با فرض برابری پارامترهای دیگر حذف می‌شود.
در قرارداد فرض‌شده سؤال، تغییر آستانه با دما نیز صفر گرفته شده است.

بدون گردکردن میانی \(kT/q\)، نسبت‌های دمای
\(\qty{300}{\kelvin}\) و \(\qty{393.15}{\kelvin}\) به‌ترتیب
حدود \(15.85\) و \(8.23\) هستند.
اختلاف کوچک با جواب‌های \(15.6\) و \(8.2\) از تقریب ولتاژ
حرارتی ناشی می‌شود. نتایج با تقریب‌های مشخص در جدول
\ref{tab:leakage-ex10-results} آمده‌اند.
\begin{table}[!htbp]
\centering
\small
\caption{نسبت جریان جدید به قدیم در تمرین \secnum{۲.۱۰}. مقایسه هر ردیف در
شرایط و دمای همان ردیف انجام شده است.}
\label{tab:leakage-ex10-results}
\begin{tabularx}{\linewidth}{>{\raggedleft\arraybackslash}Xcc}
\toprule
کمیت و شرایط & مقدار استفاده‌شده & نسبت\\
\midrule
جریان اشباع ایده‌آل & \(V_{DD}=\qty{1.2}{\volt}\) & \num{1.27}\\
نشتی زیرآستانه در اتاق & \(v_T=\qty{26}{\milli\volt}\) & \num{15.6}\\
نشتی زیرآستانه در \(\qty{120}{\degreeCelsius}\) & \(v_T=\qty{34}{\milli\volt}\) & \num{8.2}\\
\bottomrule
\end{tabularx}
\end{table}

حساسیت نسبت نشتی به کاهش آستانه با ضریب \(1/(nv_T)\) تعیین
می‌شود. در دمای بالاتر، این ضریب کوچک‌تر و نسبت تغییر ناشی از
همان کاهش آستانه کمتر است. محاسبه نشتی مطلق در دو دمای مختلف به
وابستگی دمایی \(I_{ds0}\) و دیگر پارامترهای افزاره نیز نیاز دارد.
در تغذیه ثابت و با مشخص بودن سهم جریان از منبع، توان همان مؤلفه
نشتی با همان نسبت جریان تغییر می‌کند.

\section*{مرور مدل‌ها و کنترل محاسبه}
در مدل زیرآستانه، بایاس‌ها، آستانه مؤثر، \(n\)، ولتاژ حرارتی و عامل
درین را مشخص کنید. \(S\) از \eqref{eq:leakage-subthreshold-slope}
به دست می‌آید. در نسبت جریان‌ها، عامل‌های مشترک در شرایط ثابت حذف
می‌شوند. برای \(I_{\mathrm{off}}\)، بایاس مرجع و نسبت دو عامل \(F\)
را مطابق \eqref{eq:leakage-off-normalized} نگه دارید.

در مدل گیت، ضخامت و ضرایب ماده باید واحدهای سازگار داشته باشند.
جریان کل نمودار با کمیت نرمال‌ساز محور به دست می‌آید. در مدل پیوند،
آند، کاتد و علامت ولتاژ را تعیین کنید. مدل دیود، \lr{BTBT} و
\lr{GIDL} به بایاس و پارامترهای مربوط به خود وابسته‌اند.

برای توان کل، جریان هر منبع را از \lr{KCL} محاسبه کنید.
مقادیر هر شکل با دما، فناوری و بایاس همان شکل تفسیر می‌شوند.
برای تغییر دمای همان افزاره، رابطه \eqref{eq:temperature-leakage-ratio}
و برای تغییر هندسه، بخش \ref{sec:geometry-dependence} را به کار ببرید.
